
\documentclass{article}
\usepackage{colortbl}
\usepackage{makecell}
\usepackage{multirow}
\usepackage{supertabular}

\begin{document}

\newcounter{utterance}

\centering \large Interaction Transcript for game `cladder', experiment `full\_v1.5\_default', episode 3258 with Dialogue Architects\_\_qwen3.5{-}9b\_\_SCoRe\_Qwen3.5{-}9B.
\vspace{24pt}

{ \footnotesize  \setcounter{utterance}{1}
\setlength{\tabcolsep}{0pt}
\begin{supertabular}{c@{$\;$}|p{.15\linewidth}@{}p{.15\linewidth}p{.15\linewidth}p{.15\linewidth}p{.15\linewidth}p{.15\linewidth}}
    \# & \multicolumn{2}{c}{Player} && \multicolumn{2}{c}{Game Master} \\
    \hline

    \theutterance \stepcounter{utterance}  
    & & & \multicolumn{4}{p{0.6\linewidth}}{
        \cellcolor[rgb]{0.9,0.9,0.9}{
            \makecell[{{p{\linewidth}}}]{
                \texttt{\tiny{[P1$\langle$GM]}}
                \texttt{You are an expert in causal inference. The following question is not a typical commonsense query, but rather a meticulously designed question created by a professor specializing in causal inference, intended to assess the students' mastery of the course content.} \\
\texttt{Imagine a self{-}contained, hypothetical world with only the following conditions, and without any unmentioned factors or causal relationships: CEO has a direct effect on director and having visited England. Having visited England has a direct effect on employee. Director has a direct effect on employee. The overall probability of having visited England is 19\%. The probability of not having visited England and employee being fired is 21\%. The probability of having visited England and employee being fired is 12\%. Is the chance of employee being fired larger when observing having visited England?} \\
\texttt{Start your answer with 'Yes' or 'No', followed by additional reasoning or evidence to support your explanation.} \\
            }
        }
    }
     \\ \\

    \theutterance \stepcounter{utterance}  
    & \multicolumn{4}{p{0.6\linewidth}}{
        \cellcolor[rgb]{0.561,0.737,0.561}{
            \makecell[{{p{\linewidth}}}]{
                \texttt{\tiny{[P1$\rangle$GM]}}
                \texttt{No} \\
\\ 
\texttt{**Reasoning:**} \\
\\ 
\texttt{To determine if the chance of an employee being fired is larger when observing that they have visited England, we must compare the conditional probability of being fired given a visit to England, \$P(\textbackslash{}text\{Fired\} \textbackslash{}mid \textbackslash{}text\{Visited\})\$, against the baseline probability of being fired, \$P(\textbackslash{}text\{Fired\})\$.} \\
\\ 
\texttt{Based on the provided data:} \\
\texttt{1.  **Calculate the total probability of being fired (\$P(\textbackslash{}text\{Fired\})\$):**} \\
\texttt{    The problem states the joint probabilities for the two mutually exclusive states of "having visited England" (Visited vs. Not Visited).} \\
\texttt{    *   \$P(\textbackslash{}text\{Visited\} \textbackslash{}cap \textbackslash{}text\{Fired\}) = 12\textbackslash{}\%\$} \\
\texttt{    *   \$P(\textbackslash{}text\{Not Visited\} \textbackslash{}cap \textbackslash{}text\{Fired\}) = 21\textbackslash{}\%\$} \\
\\ 
\texttt{    Therefore, the total probability of being fired is the sum of these joint probabilities:} \\
\texttt{    \$\$P(\textbackslash{}text\{Fired\}) = 12\textbackslash{}\% + 21\textbackslash{}\% = 33\textbackslash{}\%\$\$} \\
\\ 
\texttt{2.  **Calculate the conditional probability of being fired given a visit (\$P(\textbackslash{}text\{Fired\} \textbackslash{}mid \textbackslash{}text\{Visited\})\$):**} \\
\texttt{    We are given that the overall probability of having visited England is \$P(\textbackslash{}text\{Visited\}) = 19\textbackslash{}\%\$.} \\
\texttt{    Using the definition of conditional probability:} \\
\texttt{    \$\$P(\textbackslash{}text\{Fired\} \textbackslash{}mid \textbackslash{}text\{Visited\}) = \textbackslash{}frac\{P(\textbackslash{}text\{Visited\} \textbackslash{}cap \textbackslash{}text\{Fired\})\}\{P(\textbackslash{}text\{Visited\})\}\$\$} \\
\texttt{    \$\$P(\textbackslash{}text\{Fired\} \textbackslash{}mid \textbackslash{}text\{Visited\}) = \textbackslash{}frac\{12\textbackslash{}\%\}\{19\textbackslash{}\%\} \textbackslash{}approx 0.6316 \textbackslash{}text\{ or \} 63.16\textbackslash{}\%\$\$} \\
\\ 
\texttt{3.  **Comparison:**} \\
\texttt{    *   \$P(\textbackslash{}text\{Fired\} \textbackslash{}mid \textbackslash{}text\{Visited\}) \textbackslash{}approx 63.16\textbackslash{}\%\$} \\
\texttt{    *   \$P(\textbackslash{}text\{Fired\}) = 33\textbackslash{}\%\$} \\
\\ 
\texttt{Since \$63.16\textbackslash{}\% > 33\textbackslash{}\%\$, the probability of an employee being fired is indeed larger when observing that they have visited England.} \\
\\ 
\texttt{*(Note: While the causal graph describes a mechanism where CEO affects both Visited and Fired, and Visited affects Fired, the question specifically asks about the statistical association observed in the data provided. The calculation confirms a positive association between visiting England and being fired in this specific hypothetical world.)*} \\
            }
        }
    }
    & & \\ \\

    \theutterance \stepcounter{utterance}  
    & & & \multicolumn{2}{p{0.3\linewidth}}{
        \cellcolor[rgb]{0.9,0.9,0.9}{
            \makecell[{{p{\linewidth}}}]{
                \texttt{\tiny{[GM$|$GM]}}
                \texttt{No} \\
            }
        }
    }
    & & \\ \\

    \theutterance \stepcounter{utterance}  
    & & & \multicolumn{2}{p{0.3\linewidth}}{
        \cellcolor[rgb]{0.9,0.9,0.9}{
            \makecell[{{p{\linewidth}}}]{
                \texttt{\tiny{[GM$|$GM]}}
                \texttt{yes} \\
            }
        }
    }
    & & \\ \\

    \theutterance \stepcounter{utterance}  
    & & & \multicolumn{2}{p{0.3\linewidth}}{
        \cellcolor[rgb]{0.9,0.9,0.9}{
            \makecell[{{p{\linewidth}}}]{
                \texttt{\tiny{[GM$|$GM]}}
                \texttt{game\_result = LOSE} \\
            }
        }
    }
    & & \\ \\

\end{supertabular}
}

\end{document}
